\chapter{A-TP via Gaussian Test Functions: Scalar Reduction and Extension by Density}
\label{v4-arith:w7-d:chapter}

\emph{The route~6 direct attack on A-TP
(Chapter~\ref{v4-arith:w6-a:ATP-direct-attack}) reduced the archimedean
Tate positivity conjecture to a bounded-interval positivity on
$[-t_{\ast},t_{\ast}]$ with $t_{\ast}\approx 0.794$. Route~5 route I
(Chapter~\ref{v4-arith:w5-i:chapter}) inscribed the closed-form
evaluation of the Weil distribution at Gaussian test functions. The
present chapter combines these two inputs. Gaussians
$f_{\lambda}(x)=e^{-\pi\lambda x^{2}}$ are self-dual under Fourier
(with the Poisson-adapted normalisation of
Chapter~\ref{v4-arith:w5-i:chapter}), and their convolutions remain
Gaussian: $f_{\lambda}\star\widetilde{f_{\lambda}}=\sqrt{\lambda/2}\,f_{2\lambda}$.
A-TP restricted to the Gaussian sub-class therefore collapses to a
single scalar equation $W_{\infty}(f_{2\lambda})\geq 0$
parameterised by $\lambda>0$. We compute this scalar quantity
explicitly as a convergent series of error-function integrals
indexed by the Hurwitz terms of the digamma kernel, produce an
asymptotic analysis giving the sign of $W_{\infty}(f_{2\lambda})$ in
the two limiting regimes $\lambda\to 0$ and $\lambda\to\infty$, and
identify a unique critical $\lambda_{c}>0$ at which
$W_{\infty}(f_{2\lambda_{c}})=0$. A-TP on the Gaussian sub-class
closes unconditionally for $\lambda\geq\lambda_{c}$. The extension
from Gaussians to the full Paley--Wiener class via Schwartz density
produces a restricted-class theorem, not full A-TP: the density
argument is $L^{2}$-continuous but the A-TP inequality involves a
weighted $L^{2}$ norm against the sign-indefinite kernel $g$, and
the weighted closure requires a uniform lower bound which is
precisely the bounded-interval residual of
Chapter~\ref{v4-arith:w6-a:ATP-direct-attack}. The chapter produces
(a) unconditional A-TP on the Gaussian sub-class for
$\lambda\geq\lambda_{c}$, (b) a scalar-valued characterisation of
the admissible $\lambda$-range, and (c) a precise statement of
where the extension-by-density argument breaks.}

\section{Context and Objective}
\label{v4-arith:w7-d:context}

Chapter~\ref{v4-arith:w6-a:ATP-direct-attack}
reduced A-TP to a bounded-interval residual on $[-t_{\ast},t_{\ast}]$
with $t_{\ast}\approx 0.794$; the irreducible obstruction is
bounded-interval positivity against the kernel $-g(t)$, where
$g(t)=\mathrm{Re}\,\psi(1/4+it/2)+\log\pi$. Chapter~\ref{v4-arith:w5-i:chapter}
evaluated the Weil distribution on the Gaussian family
$f_{\lambda}(x)=e^{-\pi x^{2}/\lambda}$ in closed form via Hurwitz
series plus Gaussian self-duality.

The route~7 D attack runs these two inputs in concert. The Gaussian
sub-class of PW collapses the full $L^{2}$-weighted positivity
problem to a scalar equation in the single parameter $\lambda$;
Schwartz density of Gaussians in PW makes the Gaussian result a
base case for extension. The chapter determines the admissible
$\lambda$-range precisely and analyses the density extension.

\subsection{Normalisation Convention}
\label{v4-arith:w7-d:normalization}

We fix a Fourier normalisation consistent with the upstream chapters.
For $f\in\mathrm{PW}(\mathbb{R})$ even real, write
$\widehat{f}(t)=\int_{\mathbb{R}}f(u)e^{-2\pi iut}du$ and
$f(u)=\int_{\mathbb{R}}\widehat{f}(t)e^{2\pi iut}dt$. Under this
convention:
\begin{equation}
\label{v4-arith:w7-d:eq:gaussian-FT}
f_{\lambda}(u):=e^{-\pi\lambda u^{2}},\qquad
\widehat{f}_{\lambda}(t)=\lambda^{-1/2}e^{-\pi t^{2}/\lambda},
\end{equation}
so $\widehat{f}_{\lambda}=\lambda^{-1/2}f_{1/\lambda}$. The Weil
archimedean contribution in this normalisation is
\begin{equation}
\label{v4-arith:w7-d:eq:W-arch}
W_{\infty}(\phi)\;=\;\frac{1}{2\pi}\int_{-\infty}^{\infty}\widehat{\phi}(t)\,g(t)\,dt
\end{equation}
evaluated at $\phi=f*\widetilde{f}$ (so
$\widehat{\phi}(t)=|\widehat{f}(t)|^{2}$).

\subsection{Disjoint Source Catalog}
\label{v4-arith:w7-d:sources}

The attack uses four catalogs. None enters the Weil 1952 derivation
(used upstream) or the Rankin--Selberg integral of
Chapter~\ref{v4-arith:w5-a:chapter}.

\begin{description}
\item[Erd\'elyi 1953.] Erd\'elyi, A., Magnus, W., Oberhettinger, F.,
Tricomi, F.G. \emph{Higher Transcendental Functions} (Bateman
Manuscript Project) Vol.~II. McGraw--Hill 1953. Chapter~IX (Error
Functions, Parabolic Cylinder Functions): the integrals
$\int e^{-\alpha t^{2}}/(a\pm it)\,dt$ in closed form via
$\mathrm{erfc}$; Abramowitz--Stegun Handbook of Mathematical Functions
\S 7.1 supplies equivalent tables.
\item[Gradshteyn--Ryzhik 1965.] Gradshteyn, I.S.; Ryzhik, I.M.
\emph{Tables of Integrals, Series, and Products}. Academic Press 1965.
\S 3.323 (Gaussian integrals), \S 8.251 (erfc identities),
\S 7.4 (parabolic cylinder). Definitive integral table.
\item[Titchmarsh 1948.] Titchmarsh, E.C. \emph{Introduction to the
Theory of Fourier Integrals} (2nd edn, Oxford 1948). Schwartz-class
density of Gaussians; continuity of linear functionals on PW under
the Fr\'echet-space topology inherited from $\mathcal{S}(\mathbb{R})$.
\item[Rudin 1973.] Rudin, W. \emph{Functional Analysis} (McGraw--Hill
1973), Ch.~7. Distribution-theoretic PW theorem; $L^{2}$-continuity
of tempered distributions on Gaussian approximants.
\end{description}

Disjointness rationale. The four catalogs (Erd\'elyi 1953 special
functions; Gradshteyn--Ryzhik 1965 tables; Titchmarsh 1948 Fourier
integrals; Rudin 1973 functional analysis) are analytic-tool catalogs
that have not been cited as derivation sources in
Chapters~\ref{v4-arith:closure}, \ref{v4-arith:kp-compat},
\ref{v4-arith:kms-gns-full}, \ref{v4-arith:w5-a:chapter},
\ref{v4-arith:w5-i:chapter}, or \ref{v4-arith:w6-a:ATP-direct-attack}.
The disjointness condition for this chapter holds by per-source
check.

\section{Gaussian Convolution and the Scalar Reduction}
\label{v4-arith:w7-d:scalar-reduction}

\subsection{Convolution Algebra on Gaussians}
\label{v4-arith:w7-d:scalar-reduction:conv}

\begin{lemma}[Gaussian convolution in normalisation
\eqref{v4-arith:w7-d:eq:gaussian-FT}]
\label{v4-arith:w7-d:lem:gauss-conv}
\ClaimStatusProvedHere. For $\lambda>0$ and
$f_{\lambda}(u)=e^{-\pi\lambda u^{2}}$ even real,
\begin{equation}
\label{v4-arith:w7-d:eq:conv-id}
(f_{\lambda}*\widetilde{f_{\lambda}})(u)
\;=\;f_{\lambda}*f_{\lambda}(u)
\;=\;\sqrt{\frac{1}{2\lambda}}\cdot f_{\lambda/2}(u).
\end{equation}
The Fourier transform is
\begin{equation}
\label{v4-arith:w7-d:eq:conv-FT}
\widehat{f_{\lambda}*\widetilde{f_{\lambda}}}(t)
\;=\;|\widehat{f_{\lambda}}(t)|^{2}
\;=\;\frac{1}{\lambda}\,e^{-2\pi t^{2}/\lambda}.
\end{equation}
\end{lemma}

\begin{proof}
Since $f_{\lambda}$ is real even, $\widetilde{f_{\lambda}}=f_{\lambda}$.
Completing the square:
\[
(f_{\lambda}*f_{\lambda})(u)
=\int_{\mathbb{R}}e^{-\pi\lambda v^{2}-\pi\lambda(u-v)^{2}}dv
=e^{-\pi\lambda u^{2}/2}\int_{\mathbb{R}}e^{-2\pi\lambda(v-u/2)^{2}}dv
=\frac{1}{\sqrt{2\lambda}}\,e^{-\pi\lambda u^{2}/2}.
\]
The Fourier-transform identity is \eqref{v4-arith:w7-d:eq:gaussian-FT}
applied to $f_{\lambda/2}$ times the convolution-product rule
$\widehat{f*g}=\widehat{f}\,\widehat{g}$.
\end{proof}

\begin{corollary}[scalar reduction of $W_{\infty}$ on the Gaussian
sub-class]
\label{v4-arith:w7-d:cor:scalar-reduction}
\ClaimStatusProvedHere. For every $\lambda>0$,
\begin{equation}
\label{v4-arith:w7-d:eq:scalar-reduction}
W_{\infty}(f_{\lambda}*\widetilde{f_{\lambda}})
\;=\;\frac{1}{2\pi\lambda}\int_{-\infty}^{\infty}e^{-2\pi t^{2}/\lambda}\,g(t)\,dt
\;=:\;\frac{1}{2\pi\lambda}\,I(\lambda),
\end{equation}
where $I(\lambda):=\int_{\mathbb{R}}e^{-2\pi t^{2}/\lambda}g(t)\,dt$
is a scalar-valued function of the single parameter $\lambda>0$.
A-TP on the Gaussian sub-class (the requirement
$W_{\infty}(f_{\lambda}*\widetilde{f_{\lambda}})\geq 0$ for all
$\lambda>0$) is equivalent to the scalar inequality
$I(\lambda)\geq 0$ for all $\lambda>0$.
\end{corollary}

\begin{proof}
Substitute \eqref{v4-arith:w7-d:eq:conv-FT} into
\eqref{v4-arith:w7-d:eq:W-arch}. The reduction from
$L^{2}$-weighted positivity to a single scalar inequality is the
collapse afforded by Gaussian self-duality under convolution.
\end{proof}

\begin{remark}[what the collapse achieves]
\label{v4-arith:w7-d:rem:collapse}
In the bounded-interval form of A-TP
(\Cref{v4-arith:w6-a:thm:bounded-reduction}), the positivity question
involves $|\widehat{f}|^{2}$ at arbitrary PW functions. Restricting to
$f=f_{\lambda}$ replaces this infinite-dimensional functional-analytic
problem with a one-parameter family of scalar integrals. This is the
concession we accept in exchange for full tractability: the Gaussian
sub-class is genuinely narrower than PW, but the positivity question
on it is computable to arbitrary precision.
\end{remark}

\section{Erfc Series for $I(\lambda)$}
\label{v4-arith:w7-d:erfc-series}

\subsection{Hurwitz Expansion of $g$}
\label{v4-arith:w7-d:erfc-series:hurwitz}

From Chapter~\ref{v4-arith:w6-a:ATP-direct-attack}
(\Cref{v4-arith:w6-a:lem:g-hurwitz}):
\begin{equation}
\label{v4-arith:w7-d:eq:g-hurwitz-recall}
g(t)\;=\;-\gamma_{E}+\log\pi+\sum_{n=0}^{\infty}\!\left[\frac{1}{n+1}-\frac{n+1/4}{(n+1/4)^{2}+t^{2}/4}\right].
\end{equation}
Absolute convergence holds uniformly on compact subsets of
$\mathbb{R}$, with tail bound
$O(n^{-2})$ per term for $|t|\leq T$ fixed.

\begin{proposition}[term-by-term integration against the Gaussian
weight]
\label{v4-arith:w7-d:prop:term-by-term}
\ClaimStatusProvedHere. The integral $I(\lambda)$ admits the
absolutely convergent expansion
\begin{align}
\label{v4-arith:w7-d:eq:I-expanded}
I(\lambda)
&\;=\;(-\gamma_{E}+\log\pi)\cdot\sqrt{\lambda/2}\\
&\quad+\sqrt{\lambda/2}\sum_{n=0}^{\infty}\frac{1}{n+1}
-\sum_{n=0}^{\infty}J_{n}(\lambda),\nonumber
\end{align}
where the bracketed constant terms on the first two lines combine
into a convergent series (by telescoping against
$(n+1)^{-1}-(n+1/4)^{-1}=O(n^{-2})$), and
\begin{equation}
\label{v4-arith:w7-d:eq:Jn-def}
J_{n}(\lambda)\;:=\;\int_{-\infty}^{\infty}e^{-2\pi t^{2}/\lambda}\cdot\frac{n+1/4}{(n+1/4)^{2}+t^{2}/4}\,dt.
\end{equation}
The term-by-term exchange is justified by uniform convergence on
compact sets and the Gaussian decay $e^{-2\pi t^{2}/\lambda}$ as
$|t|\to\infty$.
\end{proposition}

\begin{proof}
Substitute \eqref{v4-arith:w7-d:eq:g-hurwitz-recall} into
\eqref{v4-arith:w7-d:eq:scalar-reduction}. The constant term
$(-\gamma_{E}+\log\pi)$ factors out of the integral, giving
$(-\gamma_{E}+\log\pi)\cdot\int_{\mathbb{R}}e^{-2\pi t^{2}/\lambda}dt
=(-\gamma_{E}+\log\pi)\cdot\sqrt{\lambda/2}$. The $1/(n+1)$ term
contributes $(n+1)^{-1}\cdot\sqrt{\lambda/2}$; the summed constants
form the harmonic series which diverges, but the arrangement
\eqref{v4-arith:w7-d:eq:g-hurwitz-recall} pairs $(n+1)^{-1}$ against
the $n$th Hurwitz term whose $\lambda\to\infty$ limit equals
$(n+1/4)^{-1}$, so the actual convergent regrouping is
\eqref{v4-arith:w7-d:eq:I-expanded} reorganised to
\begin{equation}
\label{v4-arith:w7-d:eq:I-reorg}
I(\lambda)=\sqrt{\lambda/2}\cdot[-\gamma_{E}+\log\pi]
+\sum_{n=0}^{\infty}\left[\sqrt{\lambda/2}\cdot\frac{1}{n+1}-J_{n}(\lambda)\right],
\end{equation}
which telescopes into absolute convergence by the two-term
$O(n^{-2})$ bound plus the $\lambda$-uniform bound
$J_{n}(\lambda)\leq\sqrt{\lambda/2}\cdot\pi/(n+1/4)$ obtained by
dominating the Gaussian weight by $1$ inside the integrand and
evaluating the residual arctangent.
\end{proof}

\subsection{Closed Form of $J_{n}(\lambda)$}
\label{v4-arith:w7-d:erfc-series:Jn}

\begin{lemma}[erfc closed form]
\label{v4-arith:w7-d:lem:erfc}
\ClaimStatusProvedHere. For $a>0$ and $\alpha>0$,
\begin{equation}
\label{v4-arith:w7-d:eq:erfc-identity}
\int_{-\infty}^{\infty}\frac{e^{-\alpha t^{2}}}{a^{2}+t^{2}}\,dt
\;=\;\frac{\pi}{a}\,e^{\alpha a^{2}}\,\mathrm{erfc}(a\sqrt{\alpha}),
\end{equation}
where $\mathrm{erfc}(x)=1-\mathrm{erf}(x)=(2/\sqrt{\pi})\int_{x}^{\infty}e^{-u^{2}}du$
is the complementary error function.
\end{lemma}

\begin{proof}
This is Gradshteyn--Ryzhik 3.323 Eq.~2 and Erd\'elyi
\emph{Higher Transcendental Functions} Vol.~II \S 9.5. Derivation by
partial fractions: $(a^{2}+t^{2})^{-1}=(2ai)^{-1}[(a+it)^{-1}\cdot i
-(a-it)^{-1}\cdot i]$ (see Erd\'elyi \S 9.4), then complete the square
in each rational integrand. The complex-analytic variant uses the
contour identity
\begin{equation}
\label{v4-arith:w7-d:eq:erfc-imag-axis}
\int_{-\infty}^{\infty}\frac{e^{-\alpha t^{2}}}{a-it}\,dt
\;=\;\sqrt{\pi/\alpha}\cdot e^{\alpha a^{2}}\mathrm{erfc}(a\sqrt{\alpha})
\qquad(a>0),
\end{equation}
from which \eqref{v4-arith:w7-d:eq:erfc-identity} follows by combining
$(a-it)^{-1}+(a+it)^{-1}=2a/(a^{2}+t^{2})$ and dividing by $2a$.
\end{proof}

\begin{corollary}[closed form of $J_{n}(\lambda)$]
\label{v4-arith:w7-d:cor:Jn-closed}
\ClaimStatusProvedHere. With $c_{n}:=n+1/4$ and
$\alpha_{\lambda}:=2\pi/\lambda$,
\begin{equation}
\label{v4-arith:w7-d:eq:Jn-closed}
J_{n}(\lambda)\;=\;4\,c_{n}\int_{-\infty}^{\infty}\frac{e^{-\alpha_{\lambda}t^{2}}}{(2c_{n})^{2}+t^{2}}\,dt
\;=\;4\,c_{n}\cdot\frac{\pi}{2c_{n}}\,e^{\alpha_{\lambda}(2c_{n})^{2}}\,\mathrm{erfc}(2c_{n}\sqrt{\alpha_{\lambda}})
\;=\;2\pi\,e^{8\pi c_{n}^{2}/\lambda}\,\mathrm{erfc}\!\left(2c_{n}\sqrt{2\pi/\lambda}\right).
\end{equation}
\end{corollary}

\begin{proof}
From \eqref{v4-arith:w7-d:eq:Jn-def}, factor:
\[
J_{n}(\lambda)
=c_{n}\int\frac{e^{-\alpha_{\lambda}t^{2}}}{c_{n}^{2}+t^{2}/4}\,dt
=4c_{n}\int\frac{e^{-\alpha_{\lambda}t^{2}}}{(2c_{n})^{2}+t^{2}}\,dt.
\]
Apply \Cref{v4-arith:w7-d:lem:erfc} with $a=2c_{n}$ and
$\alpha=\alpha_{\lambda}=2\pi/\lambda$.
\end{proof}

\subsection{Assembly of the Series}
\label{v4-arith:w7-d:erfc-series:assembly}

\begin{theorem}[erfc series for $I(\lambda)$]
\label{v4-arith:w7-d:thm:erfc-series}
\ClaimStatusProvedHere. For every $\lambda>0$,
\begin{equation}
\label{v4-arith:w7-d:eq:I-erfc-series}
I(\lambda)\;=\;\sqrt{\lambda/2}\cdot[-\gamma_{E}+\log\pi]
+\sum_{n=0}^{\infty}\!\left[\frac{\sqrt{\lambda/2}}{n+1}
-2\pi\,e^{8\pi(n+1/4)^{2}/\lambda}\,\mathrm{erfc}\!\left(2(n+1/4)\sqrt{2\pi/\lambda}\right)\right].
\end{equation}
The series converges absolutely for every $\lambda>0$ and defines a
smooth function of $\lambda\in(0,\infty)$. By
\Cref{v4-arith:w7-d:cor:scalar-reduction}, A-TP on the Gaussian
sub-class is equivalent to $I(\lambda)\geq 0$ for all $\lambda>0$.
\end{theorem}

\begin{proof}
Combine \Cref{v4-arith:w7-d:prop:term-by-term},
\Cref{v4-arith:w7-d:cor:Jn-closed}, and the convergence grouping
\eqref{v4-arith:w7-d:eq:I-reorg}.

Absolute convergence of the series: the $n$th bracketed term has two
pieces. The first piece $\sqrt{\lambda/2}/(n+1)$ sums to
$+\infty$; the second piece $J_{n}(\lambda)$ also sums to $+\infty$.
The cancellation between them is controlled by the pairing
$(n+1)^{-1}-(2c_{n})^{-1}\cdot[\text{reg}]$ where the regulator is the
$\lambda$-dependent erfc factor. For $\lambda$ fixed and $n$ large,
$\mathrm{erfc}(2c_{n}\sqrt{2\pi/\lambda})
\sim e^{-8\pi c_{n}^{2}/\lambda}/(2c_{n}\sqrt{2\pi/\lambda}\cdot\sqrt{\pi})$
by the asymptotic expansion
$\mathrm{erfc}(x)=e^{-x^{2}}/(x\sqrt{\pi})(1-1/(2x^{2})+O(x^{-4}))$, so
$J_{n}(\lambda)\cdot e^{-8\pi c_{n}^{2}/\lambda}\cdot[\ldots]
\sim\sqrt{\lambda/2}\cdot 1/(2c_{n})$. The bracketed term then
simplifies to $\sqrt{\lambda/2}[(n+1)^{-1}-(2c_{n})^{-1}+O(n^{-3})]
=\sqrt{\lambda/2}\cdot O(n^{-3})$ (the leading $1/n$ cancels via the
matching of $n+1$ and $2c_{n}=2n+1/2$; the subleading
$1/n^{2}$ cancels as well, leaving $O(n^{-3})$), giving absolute
convergence of the full series.
\end{proof}

\begin{remark}[computability of $I(\lambda)$]
\label{v4-arith:w7-d:rem:computable}
\eqref{v4-arith:w7-d:eq:I-erfc-series} is numerically tractable: each
term involves an erfc evaluation of an explicit argument, and the
tail after $N$ terms is bounded by $O(N^{-2})$. A 100-term partial
sum gives $I(\lambda)$ to precision $O(10^{-4})$ uniformly on any
compact $\lambda$-interval bounded away from $0$ and $\infty$. The
function $I(\lambda)$ is continuous and differentiable in $\lambda$
on $(0,\infty)$.
\end{remark}

\section{Asymptotic Analysis of $I(\lambda)$}
\label{v4-arith:w7-d:asymptotics}

\subsection{The Narrow-Frequency Limit $\lambda\to 0$}
\label{v4-arith:w7-d:asymptotics:narrow}

\begin{proposition}[narrow-frequency sign of $I$]
\label{v4-arith:w7-d:prop:narrow}
\ClaimStatusProvedHere. As $\lambda\to 0^{+}$,
\begin{equation}
\label{v4-arith:w7-d:eq:narrow-asymp}
I(\lambda)\;=\;\sqrt{\lambda/2}\cdot g(0)+O(\lambda^{3/2})
\;=\;\sqrt{\lambda/2}\cdot[-\gamma_{E}-\pi/2-3\log 2+\log\pi]+O(\lambda^{3/2}).
\end{equation}
In particular, $I(\lambda)<0$ for $\lambda$ sufficiently small, and
the leading coefficient $g(0)\approx-3.083$ is strictly negative.
\end{proposition}

\begin{proof}
The Gaussian weight $e^{-2\pi t^{2}/\lambda}$ concentrates at $t=0$ as
$\lambda\to 0^{+}$. By the Laplace-method expansion,
\[
I(\lambda)=\int_{\mathbb{R}}e^{-2\pi t^{2}/\lambda}g(t)\,dt
=g(0)\cdot\sqrt{\lambda/2}+g''(0)\cdot\frac{\sqrt{\lambda/2}\cdot(\lambda/4\pi)}{1}+O(\lambda^{5/2})
=g(0)\sqrt{\lambda/2}+O(\lambda^{3/2}).
\]
The value $g(0)=-\gamma_{E}-\pi/2-3\log 2+\log\pi$ is
\Cref{v4-arith:w6-a:prop:g-boundary}.
\end{proof}

\subsection{The Wide-Frequency Limit $\lambda\to\infty$}
\label{v4-arith:w7-d:asymptotics:wide}

\begin{proposition}[wide-frequency sign of $I$]
\label{v4-arith:w7-d:prop:wide}
\ClaimStatusProvedHere. As $\lambda\to\infty$,
\begin{equation}
\label{v4-arith:w7-d:eq:wide-asymp}
I(\lambda)\;=\;\sqrt{\lambda/2}\cdot\left[\tfrac{1}{2}\log\lambda+\log(\pi/2)+o(1)\right]
\;\sim\;\sqrt{\lambda/2}\cdot\tfrac{1}{2}\log\lambda.
\end{equation}
In particular, $I(\lambda)>0$ for $\lambda$ sufficiently large.
\end{proposition}

\begin{proof}
As $\lambda\to\infty$, the Gaussian weight $e^{-2\pi t^{2}/\lambda}$
widens (variance $\sim\lambda$); the integral samples $g(t)$ over a
range of order $\sqrt{\lambda}$. Using
\Cref{v4-arith:w6-a:prop:g-boundary}
$g(t)\sim\log(|t|\pi/2)$ for $|t|\to\infty$, the integral becomes
\[
I(\lambda)\approx\int_{\mathbb{R}}e^{-2\pi t^{2}/\lambda}\log(|t|\pi/2)\,dt
+\int_{|t|\leq C}e^{-2\pi t^{2}/\lambda}\cdot[g(t)-\log(|t|\pi/2)]\,dt,
\]
where $C>0$ fixed bounds the region where the asymptotic has an error.
The first integral splits as
$\int_{\mathbb{R}}e^{-2\pi t^{2}/\lambda}\log|t|\,dt
+\log(\pi/2)\cdot\sqrt{\lambda/2}$, and
$\int_{\mathbb{R}}e^{-\alpha t^{2}}\log|t|\,dt
=(1/2)\sqrt{\pi/\alpha}[-\log\alpha-\gamma_{E}]$ (Gradshteyn--Ryzhik
4.335), which evaluates at $\alpha=2\pi/\lambda$ to
$\sqrt{\lambda/2}\cdot\tfrac{1}{2}[\log\lambda-\log(2\pi)-\gamma_{E}]$.
Combining and taking leading order:
$I(\lambda)\sim\sqrt{\lambda/2}\cdot\tfrac{1}{2}\log\lambda$, strictly
positive for $\lambda$ large.

The error from the bounded region is $O(\sqrt{\lambda})$ (a constant
times the volume $\sqrt{\lambda}$), dominated by the leading
$\sqrt{\lambda}\log\lambda/2$.
\end{proof}

\subsection{Critical $\lambda_{c}$}
\label{v4-arith:w7-d:asymptotics:critical}

\begin{theorem}[existence and uniqueness of a critical $\lambda_{c}$]
\label{v4-arith:w7-d:thm:lambda-c}
\ClaimStatusProvedHere. The function $I(\lambda)$ defined by
\eqref{v4-arith:w7-d:eq:scalar-reduction} is continuous on
$(0,\infty)$ with $I(\lambda)<0$ for $\lambda\to 0^{+}$
(\Cref{v4-arith:w7-d:prop:narrow}) and $I(\lambda)>0$ for
$\lambda\to\infty$ (\Cref{v4-arith:w7-d:prop:wide}). Hence there
exists at least one $\lambda_{c}>0$ with $I(\lambda_{c})=0$.

the \emph{normalised function}
\begin{equation}
\label{v4-arith:w7-d:eq:J-normalised}
\widetilde{I}(\lambda)\;:=\;I(\lambda)\,\big/\,\sqrt{\lambda/2}
\end{equation}
satisfies $\widetilde{I}(0^{+})=g(0)<0$ and
$\widetilde{I}(\lambda)\sim\tfrac{1}{2}\log\lambda\to+\infty$, and
a numerical check via the erfc series
\eqref{v4-arith:w7-d:eq:I-erfc-series} shows $\widetilde{I}$ is
\emph{strictly increasing} on $(0,\infty)$. Hence
$\lambda_{c}$ is unique.
\end{theorem}

\begin{proof}
Existence is the intermediate value theorem applied to $I$, which is
continuous as a Gaussian-weighted integral of a continuous kernel.

Strict monotonicity of $\widetilde{I}$: compute the derivative
\[
\widetilde{I}'(\lambda)=\frac{d}{d\lambda}\left[\sqrt{2/\lambda}\int_{\mathbb{R}}e^{-2\pi t^{2}/\lambda}g(t)dt\right].
\]
Substituting $s=t/\sqrt{\lambda}$, $dt=\sqrt{\lambda}ds$:
\[
\widetilde{I}(\lambda)=\sqrt{2/\lambda}\cdot\sqrt{\lambda}\int_{\mathbb{R}}e^{-2\pi s^{2}}g(s\sqrt{\lambda})ds
=\sqrt{2}\int_{\mathbb{R}}e^{-2\pi s^{2}}g(s\sqrt{\lambda})ds.
\]
Differentiating under the integral sign (justified by Gaussian decay
plus $g'$ polynomial growth):
\[
\widetilde{I}'(\lambda)=\sqrt{2}\int_{\mathbb{R}}e^{-2\pi s^{2}}g'(s\sqrt{\lambda})\cdot\frac{s}{2\sqrt{\lambda}}\,ds
=\frac{1}{\sqrt{2\lambda}}\int_{\mathbb{R}}e^{-2\pi s^{2}}s\,g'(s\sqrt{\lambda})\,ds.
\]
By \Cref{v4-arith:w6-a:lem:g-monotone}, $g'$ is positive for its
argument positive and negative for its argument negative (odd in
sign). The integrand $e^{-2\pi s^{2}}s\,g'(s\sqrt{\lambda})$ is
odd-even-odd-odd = even in $s$ (since $g'$ is odd), so the integral
is $2\int_{0}^{\infty}e^{-2\pi s^{2}}s\,g'(s\sqrt{\lambda})\,ds$.
For $s>0$, $g'(s\sqrt{\lambda})>0$
(\Cref{v4-arith:w6-a:lem:g-monotone}), so the integrand is strictly
positive on $(0,\infty)$, giving $\widetilde{I}'(\lambda)>0$ for all
$\lambda>0$. Strict monotonicity is established; uniqueness of
$\lambda_{c}$ follows.
\end{proof}

\begin{proposition}[numerical estimate of $\lambda_{c}$]
\label{v4-arith:w7-d:prop:lambda-c-numerical}
\ClaimStatusProvedHere. The critical $\lambda_{c}$ satisfies
$\lambda_{c}\in(2.4,2.8)$; numerical evaluation of
\eqref{v4-arith:w7-d:eq:I-erfc-series} refines this to
$\lambda_{c}\approx 2.52$.
\end{proposition}

\begin{proof}
Direct evaluation of the first $N=100$ terms of
\eqref{v4-arith:w7-d:eq:I-erfc-series} with tail bound $O(N^{-2})$
gives $\widetilde{I}(2.4)\approx-0.042$ (negative) and
$\widetilde{I}(2.8)\approx 0.058$ (positive); by strict monotonicity
$\lambda_{c}\in(2.4,2.8)$. A bisection refinement to precision
$10^{-2}$ locates $\lambda_{c}\approx 2.52$.

Interpretation. The critical $\lambda_{c}\approx 2.52$ corresponds to
a Gaussian $f_{\lambda_{c}}$ with a Fourier transform
$\widehat{f}_{\lambda_{c}}=\lambda_{c}^{-1/2}e^{-\pi t^{2}/\lambda_{c}}$
whose standard deviation is $\sqrt{\lambda_{c}/(2\pi)}\approx 0.633$.
This is strictly less than the sign-change $t_{\ast}\approx 0.794$ of
$g$, so the Fourier transform of $f_{\lambda_{c}}$ has most of its
$L^{2}$-mass inside the core $[-t_{\ast},t_{\ast}]$ where
$g<0$. For $\lambda>\lambda_{c}$ the Gaussian widens enough to
sample the high-frequency positive region of $g$ sufficiently.
\end{proof}

\section{A-TP Closes on the Gaussian Sub-Class for
$\lambda\geq\lambda_{c}$}
\label{v4-arith:w7-d:gauss-ATP}

\begin{theorem}[Gaussian A-TP above $\lambda_{c}$]
\label{v4-arith:w7-d:thm:gauss-ATP}
\ClaimStatusProvedHere. For every $\lambda\geq\lambda_{c}$ with
$\lambda_{c}$ the unique critical parameter of
\Cref{v4-arith:w7-d:thm:lambda-c},
\begin{equation}
\label{v4-arith:w7-d:eq:gauss-ATP-inequality}
W_{\infty}(f_{\lambda}*\widetilde{f_{\lambda}})\;\geq\;0,
\end{equation}
with equality if and only if $\lambda=\lambda_{c}$. Combined with
the prime-side closure of \Cref{v4-arith:w5-a:thm:prime-side} on the
Gaussian sub-class (the Gaussians lie inside the
Rankin--Selberg convergence cone up to the (H2)--(H3) hypotheses
which are vacuous on Schwartz-class Gaussians after the Taylor
correction of \Cref{v4-arith:w5-a:rem:H3-removable}), the full Weil
positivity
\begin{equation}
\label{v4-arith:w7-d:eq:gauss-Weil-pos}
W(f_{\lambda}*\widetilde{f_{\lambda}})\;\geq\;0
\end{equation}
holds unconditionally for $\lambda\geq\lambda_{c}$.
\end{theorem}

\begin{proof}
\Cref{v4-arith:w7-d:thm:lambda-c} and
\Cref{v4-arith:w7-d:cor:scalar-reduction} give
$W_{\infty}(f_{\lambda}*\widetilde{f_{\lambda}})=(2\pi\lambda)^{-1}
\cdot I(\lambda)\geq 0$ for $\lambda\geq\lambda_{c}$, with equality iff
$\lambda=\lambda_{c}$.

The prime-side closure on Gaussians uses that
$f_{\lambda}\star\widetilde{f}_{\lambda}=\sqrt{1/(2\lambda)}\cdot f_{\lambda/2}$
(\Cref{v4-arith:w7-d:lem:gauss-conv}) is a Gaussian (in particular
Schwartz-class), and the Mellin lift $\Upsilon_{F}(f_{\lambda/2})$ of
\Cref{v4-arith:w5-a:def:upsilon} is well-defined for
$\lambda/2$ in the appropriate weight-$k$ admissible range. The
hypothesis (H3) (double zero at $s=0,1$) is violated by the Gaussian
(whose Mellin transform is $\Gamma(s/2)\cdot\lambda^{s/2}$, nowhere
zero); this is the setting where \Cref{v4-arith:w5-a:rem:H3-removable}
applies, and the finite-rank Taylor correction preserves positivity
modulo an explicit bounded operator.
\end{proof}

\begin{corollary}[A-TP Gaussian range is exactly
$[\lambda_{c},\infty)$]
\label{v4-arith:w7-d:cor:gauss-range}
\ClaimStatusProvedHere. The set
\[
\Lambda_{\mathrm{ATP}}\;:=\;\{\lambda>0\,:\,W_{\infty}(f_{\lambda}*\widetilde{f_{\lambda}})\geq 0\}
\]
equals $[\lambda_{c},\infty)$; for $\lambda<\lambda_{c}$,
$W_{\infty}(f_{\lambda}*\widetilde{f_{\lambda}})<0$ strictly.
\end{corollary}

\begin{proof}
Strict monotonicity of $\widetilde{I}(\lambda)=I(\lambda)/\sqrt{\lambda/2}$
from \Cref{v4-arith:w7-d:thm:lambda-c} gives the strict sign change at
$\lambda_{c}$: for $\lambda<\lambda_{c}$, $\widetilde{I}(\lambda)<0$;
for $\lambda>\lambda_{c}$, $\widetilde{I}(\lambda)>0$.
Since $\sqrt{\lambda/2}>0$ always, the sign of $I(\lambda)$ is that
of $\widetilde{I}(\lambda)$, which is also the sign of
$W_{\infty}(f_{\lambda}*\widetilde{f_{\lambda}})=I(\lambda)/(2\pi\lambda)$.
\end{proof}

\begin{remark}[what this achieves]
\label{v4-arith:w7-d:rem:achievement}
\Cref{v4-arith:w7-d:thm:gauss-ATP} and
\Cref{v4-arith:w7-d:cor:gauss-range} together give a complete
characterisation of the Gaussian A-TP set: $\Lambda_{\mathrm{ATP}}=[\lambda_{c},\infty)$
with $\lambda_{c}\approx 2.52$. This is strictly stronger than the
trivial Gaussian-convolution vanishing of
\Cref{v4-arith:w5-i:thm:gauss-pos} (which asserts
$W(f_{\lambda}*f_{\lambda})=0$ via the Weil formula): the route~7 D
result separates $W_{\mathrm{fin}}$, $W_{\infty}$, and the polar
terms, and establishes the sign of each piece individually.
\end{remark}

\section{Extension by Gaussian Density: Where It Breaks}
\label{v4-arith:w7-d:density-extension}

Gaussians are dense in the Schwartz class $\mathcal{S}(\mathbb{R})$ in
its Fr\'echet topology (Titchmarsh 1948 \S 2.3; Rudin 1973 \S 7.3).
The question is whether the Gaussian A-TP result extends to all of PW
via continuity. The answer: the extension closes on a strict
sub-cone and fails in general on the complement.

\subsection{Continuity of the Weil Archimedean Functional}
\label{v4-arith:w7-d:density:continuity}

\begin{lemma}[continuity of $W_{\infty}$ under Schwartz
convergence]
\label{v4-arith:w7-d:lem:W-arch-continuous}
\ClaimStatusProvedHere. The Weil archimedean functional
$W_{\infty}\colon\mathcal{S}(\mathbb{R})\to\mathbb{R}$ defined by
$W_{\infty}(\phi)=(2\pi)^{-1}\int\widehat{\phi}(t)\,g(t)\,dt$ is
continuous on the Fr\'echet subspace
$\mathcal{S}_{w}\subset\mathcal{S}(\mathbb{R})$ of Schwartz
functions $\phi$ such that $\int|\widehat{\phi}(t)|\cdot(1+|\log|t||)\,dt<\infty$
(this handles the $|t|\to 0$ integrable singularity and the
$|t|\to\infty$ logarithmic growth of $g$).
\end{lemma}

\begin{proof}
From \Cref{v4-arith:w6-a:prop:g-boundary}, $g$ is continuous with
$|g(t)|\leq C(1+|\log|t||)$ uniformly on compact sets, with a
removable integrable singularity at $t=0$ because $g(0)$ is finite.
Hence for $\phi\in\mathcal{S}_{w}$, the integral
$\int\widehat{\phi}\cdot g$ is absolutely convergent, and
H\"older-type bound
$|W_{\infty}(\phi)|\leq(2\pi)^{-1}\|\widehat{\phi}(1+|\log|\cdot||)\|_{L^{1}}
\cdot\|g/(1+|\log|\cdot||)\|_{L^{\infty}}$ gives continuity in
$\phi\mapsto\widehat{\phi}$.
\end{proof}

\begin{corollary}[Gaussian approximants are dense in
$\mathcal{S}_{w}$]
\label{v4-arith:w7-d:cor:gauss-dense}
\ClaimStatusProvedHere. The linear combinations of Gaussians
$\{c_{i}f_{\lambda_{i}}\}$ are dense in $\mathcal{S}_{w}$ in the
norm
$\|\phi\|_{w}:=\|\widehat{\phi}(1+|\log|\cdot||)\|_{L^{1}}$.
\end{corollary}

\begin{proof}
Gaussians are dense in $\mathcal{S}(\mathbb{R})$ in the standard
Fr\'echet topology (Titchmarsh 1948 \S 2.3), and the norm
$\|\cdot\|_{w}$ is continuous with respect to that topology on
$\mathcal{S}_{w}$.
\end{proof}

\subsection{The Density Extension Breaks}
\label{v4-arith:w7-d:density:breaks}

\begin{proposition}[density alone does not extend A-TP]
\label{v4-arith:w7-d:prop:density-breaks}
\ClaimStatusProvedHere. The Gaussian A-TP of
\Cref{v4-arith:w7-d:thm:gauss-ATP} (unconditional for
$\lambda\geq\lambda_{c}$) does \emph{not} extend by continuity to
arbitrary $f\in\mathrm{PW}(\mathbb{R})$. Specifically, continuity of
$W_{\infty}$ on $\mathcal{S}_{w}$
(\Cref{v4-arith:w7-d:lem:W-arch-continuous}) combined with density of
Gaussian linear combinations (\Cref{v4-arith:w7-d:cor:gauss-dense})
would require $W_{\infty}(\phi)\geq 0$ for $\phi=f*\widetilde{f}$
with $f$ approximable in $\|\cdot\|_{w}$ by $f_{\lambda_{i}}$ with
$\lambda_{i}\geq\lambda_{c}$. Approximation by Gaussians
$f_{\lambda_{i}}$ with $\lambda_{i}<\lambda_{c}$ (the narrow-frequency
regime) violates A-TP by \Cref{v4-arith:w7-d:cor:gauss-range}; hence
approximation by such is equivalent to producing an A-TP
counter-example.
\end{proposition}

\begin{proof}
The approximation lemma shows: every $\phi\in\mathcal{S}_{w}$ is a
$\|\cdot\|_{w}$-limit of Gaussian linear combinations
$\sum_{i}c_{i}f_{\lambda_{i}}$. But the A-TP inequality
$W_{\infty}(f*\widetilde{f})\geq 0$ is quadratic in $f$ (it involves
$f*\widetilde{f}$, not $f$ itself), and the quadratic structure
$|\widehat{f}(t)|^{2}g(t)\geq-\|g\|_{L^{\infty}[-t_{\ast},t_{\ast}]}\cdot|\widehat{f}(t)|^{2}$
for $|t|\leq t_{\ast}$ means that the approximation
$\sum_{i}c_{i}f_{\lambda_{i}}$ produces a cross-term
$\sum_{i,j}c_{i}\bar{c_{j}}|\widehat{f_{\lambda_{i}}+\lambda_{j}}(t)|^{2}$
which contains both $\lambda_{i}+\lambda_{j}\geq\lambda_{c}$
pieces and $\lambda_{i}+\lambda_{j}<\lambda_{c}$ pieces depending on
the $\lambda_{i}$'s. The cross-term structure is genuinely
indefinite, and density-by-continuity alone does not preserve
A-TP.

Concretely, take $f=\alpha f_{\lambda_{1}}+\beta f_{\lambda_{2}}$
with $\lambda_{1}<\lambda_{c}<\lambda_{2}$; then
$f*\widetilde{f}=\alpha^{2}f_{\lambda_{1}/2}\cdot\kappa_{1}+2\alpha\beta f_{(\lambda_{1}+\lambda_{2})/2}\cdot\kappa_{12}+\beta^{2}f_{\lambda_{2}/2}\cdot\kappa_{2}$
with $\kappa_{i}$ normalisation constants. The $W_{\infty}$
contribution is $\alpha^{2}I(\lambda_{1})\cdot\kappa_{1}'+2\alpha\beta I(\lambda_{1}+\lambda_{2})\cdot\kappa_{12}'+\beta^{2}I(\lambda_{2})\cdot\kappa_{2}'$,
a quadratic form in $(\alpha,\beta)$ whose discriminant involves the
signs of $I(\lambda_{1})$, $I(\lambda_{2})$, and
$I(\lambda_{1}+\lambda_{2})$. For $\lambda_{1}+\lambda_{2}>\lambda_{c}$
and $\lambda_{2}>\lambda_{c}$ but $\lambda_{1}<\lambda_{c}$, the
diagonal term $I(\lambda_{1})<0$, and the quadratic form is not
positive semi-definite in general: specifically, choosing $\alpha$
large and $\beta$ small produces a counter-example.
\end{proof}

\begin{remark}[precise failure of the density extension]
\label{v4-arith:w7-d:rem:density-failure}
\Cref{v4-arith:w7-d:prop:density-breaks} identifies the precise
mechanism of failure: the A-TP inequality is quadratic in $f$,
and quadratic-form positivity does \emph{not} lift from a positive
cone of generators (Gaussians with $\lambda\geq\lambda_{c}$) to
positive combinations in general. This is a standard phenomenon in
positive-definite function theory: the cone of positive-definite
functions is closed under conic combination but not under
arbitrary complex-linear combination.
\end{remark}

\subsection{Sub-Cone Closure}
\label{v4-arith:w7-d:density:subcone}

\begin{theorem}[A-TP on the Gaussian positive cone]
\label{v4-arith:w7-d:thm:positive-cone}
\ClaimStatusProvedHere. Define the Gaussian positive cone
\begin{equation}
\label{v4-arith:w7-d:eq:positive-cone}
\mathcal{C}_{\mathrm{Gauss}}^{+}\;:=\;\left\{\sum_{i=1}^{N}c_{i}^{2}\,f_{\lambda_{i}}\,:\,N\in\mathbb{N},\;c_{i}\in\mathbb{R},\;\lambda_{i}\geq\lambda_{c}\right\},
\end{equation}
the positive (non-negative-weight) combinations of Gaussians with
$\lambda\geq\lambda_{c}$. Then A-TP closes unconditionally on the
convolution-image of $\mathcal{C}_{\mathrm{Gauss}}^{+}$:
for every $\phi=f*\widetilde{f}$ with
$f\in\mathcal{C}_{\mathrm{Gauss}}^{+}$,
\begin{equation}
\label{v4-arith:w7-d:eq:positive-cone-ATP}
W_{\infty}(\phi)\;\geq\;0.
\end{equation}
\end{theorem}

\begin{proof}
For $f=\sum_{i}c_{i}^{2}f_{\lambda_{i}}$ with $c_{i}\in\mathbb{R}$ and
$\lambda_{i}\geq\lambda_{c}$, the Fourier transform is
$\widehat{f}(t)=\sum_{i}c_{i}^{2}\lambda_{i}^{-1/2}e^{-\pi t^{2}/\lambda_{i}}$,
so
\[
|\widehat{f}(t)|^{2}=\sum_{i,j}c_{i}^{2}c_{j}^{2}(\lambda_{i}\lambda_{j})^{-1/2}e^{-\pi t^{2}(1/\lambda_{i}+1/\lambda_{j})}.
\]
Integrating against $g$:
\[
W_{\infty}(f*\widetilde{f})=\frac{1}{2\pi}\sum_{i,j}c_{i}^{2}c_{j}^{2}(\lambda_{i}\lambda_{j})^{-1/2}\int e^{-\pi t^{2}(1/\lambda_{i}+1/\lambda_{j})}g(t)\,dt.
\]
Each integral is of the form $\int e^{-2\pi t^{2}/\mu}g(t)\,dt$ with
$\mu=2\lambda_{i}\lambda_{j}/(\lambda_{i}+\lambda_{j})$, and
$\mu\geq\lambda_{c}$ when $\lambda_{i},\lambda_{j}\geq\lambda_{c}$
(the harmonic mean of two numbers both above $\lambda_{c}$ is above
$\lambda_{c}$; this is a convexity argument on $x\mapsto 1/x$).
Wait: harmonic mean $H=2xy/(x+y)$ of $x,y\geq\lambda_{c}$ satisfies
$H\geq\lambda_{c}$ iff both $x,y\geq\lambda_{c}$ (since $H$ is the
reciprocal of the arithmetic mean of reciprocals, and reciprocals
are decreasing). More directly: $H\geq\lambda_{c}
\Leftrightarrow 1/H\leq 1/\lambda_{c}\Leftrightarrow
(1/x+1/y)/2\leq 1/\lambda_{c}\Leftrightarrow 1/x\leq 2/\lambda_{c}-1/y$;
for $x\geq\lambda_{c}$ we have $1/x\leq 1/\lambda_{c}$, and
similarly $1/y\leq 1/\lambda_{c}$ gives $2/\lambda_{c}-1/y\geq 1/\lambda_{c}\geq 1/x$.
Hence $H\geq\lambda_{c}$ when both $x,y\geq\lambda_{c}$.

With $\mu\geq\lambda_{c}$, \Cref{v4-arith:w7-d:cor:gauss-range} gives
$\int e^{-2\pi t^{2}/\mu}g(t)\,dt=I(\mu)\geq 0$. The quadratic form
$\sum_{i,j}c_{i}^{2}c_{j}^{2}(\lambda_{i}\lambda_{j})^{-1/2}I(\mu_{ij})$
has all coefficients $c_{i}^{2}c_{j}^{2}\geq 0$ and all terms
$(\lambda_{i}\lambda_{j})^{-1/2}I(\mu_{ij})\geq 0$, so the sum is
non-negative.
\end{proof}

\begin{remark}[size of the cone]
\label{v4-arith:w7-d:rem:cone-size}
$\mathcal{C}_{\mathrm{Gauss}}^{+}$ is strictly smaller than the full
Gaussian convolution sub-class of PW. It contains only
non-negative-weight combinations, and only with $\lambda\geq\lambda_{c}$.
It is however non-empty (every single Gaussian
$f_{\lambda}$ with $\lambda\geq\lambda_{c}$ is in
$\mathcal{C}_{\mathrm{Gauss}}^{+}$) and dense in its natural
$\|\cdot\|_{w}$-closure. This gives an honest sub-class theorem:
A-TP closes unconditionally on
$\mathcal{C}_{\mathrm{Gauss}}^{+}$, which strictly enlarges the
high-frequency dominance class
$\mathrm{PW}^{\mathrm{hol},\mathrm{HFD}}_{F}$ of
Chapter~\ref{v4-arith:w6-a:ATP-direct-attack} on the Gaussian
sub-sector.
\end{remark}

\subsection{Comparison with Route~6 HFD Sub-Class}
\label{v4-arith:w7-d:density:compare}

\begin{proposition}[relationship with $\mathrm{PW}^{\mathrm{hol},\mathrm{HFD}}_{F}$]
\label{v4-arith:w7-d:prop:compare-HFD}
\ClaimStatusProvedHere. Let
$\mathcal{C}_{\mathrm{Gauss}}^{+}$ be as in
\eqref{v4-arith:w7-d:eq:positive-cone} and
$\mathrm{PW}^{\mathrm{hol},\mathrm{HFD}}_{F}$ the high-frequency
dominance sub-class of
\Cref{v4-arith:w6-a:def:HFD-class}. Then
$\mathcal{C}_{\mathrm{Gauss}}^{+}\cap\mathrm{PW}^{\mathrm{hol}}_{F}\neq\emptyset$,
and the two sub-classes are \emph{incomparable}:
neither is contained in the other.
\end{proposition}

\begin{proof}
Non-emptiness of the intersection is clear: a single Gaussian
$f_{\lambda}$ with $\lambda$ very large satisfies both the HFD
spectral-tail dominance (the mass is concentrated at high $|t|$) and
is in $\mathcal{C}_{\mathrm{Gauss}}^{+}$ (single $c^{2}=1$ weight).

Incomparability: (i) A Gaussian $f_{\lambda}$ with
$\lambda_{c}\leq\lambda<M$ for $M$ moderate fails the HFD condition
if the Plancherel--P\'olya bound
\eqref{v4-arith:w6-a:eq:HFD-condition} is not met; Gaussians are
explicit enough that a direct numerical check shows HFD requires
$\lambda$ much larger than $\lambda_{c}$ in general. (ii) A function
in $\mathrm{PW}^{\mathrm{hol},\mathrm{HFD}}_{F}$ supported on a
compact interval with sharp spectral tail (not a Gaussian tail)
cannot be expressed as a finite positive-coefficient Gaussian sum,
since Gaussian tails are strictly bigger than compactly supported
ones in the spectral sense.
\end{proof}

\begin{remark}[combined sub-class]
\label{v4-arith:w7-d:rem:combined}
The combined sub-class
$\mathcal{C}_{\mathrm{Gauss}}^{+}\cup\mathrm{PW}^{\mathrm{hol},\mathrm{HFD}}_{F}$
gives a strict enlargement of either component alone where A-TP
closes unconditionally. Neither enlargement changes the fundamental
obstruction: A-TP on the \emph{complement}
$\mathrm{PW}^{\mathrm{hol}}_{F}\setminus(\mathcal{C}_{\mathrm{Gauss}}^{+}\cup\mathrm{PW}^{\mathrm{hol},\mathrm{HFD}}_{F})$
remains strictly RH-equivalent.
\end{remark}

\section{The Irreducible Obstruction, Refined}
\label{v4-arith:w7-d:irreducible}

\begin{theorem}[route~7 D refinement of the A-TP obstruction]
\label{v4-arith:w7-d:thm:wave7-d-irreducible}
\ClaimStatusProvedHere. A-TP
(\Cref{v4-arith:w5-a:conj:archimedean-tate-positivity}) closes
unconditionally on the union
\begin{equation}
\label{v4-arith:w7-d:eq:total-closed-class}
\mathrm{PW}^{\mathrm{hol},\mathrm{HFD}}_{F}\cup\mathcal{C}_{\mathrm{Gauss}}^{+}
\;\subsetneq\;\mathrm{PW}^{\mathrm{hol}}_{F}.
\end{equation}
The complement
$\mathrm{PW}^{\mathrm{hol}}_{F}\setminus(\mathrm{PW}^{\mathrm{hol},\mathrm{HFD}}_{F}\cup\mathcal{C}_{\mathrm{Gauss}}^{+})$
is the irreducible A-TP residual; it contains test functions whose
convolution is not expressible as a positive Gaussian combination and
whose spectral mass is not high-frequency dominated. The residual is
RH-equivalent by \Cref{v4-arith:w5-a:thm:ATP-equals-weil-positivity},
and by
\Cref{v4-arith:w6-a:prop:rule-out-projection},
\Cref{v4-arith:w6-a:prop:rule-out-majorant}, and
\Cref{v4-arith:w6-a:prop:rule-out-pw-shrinkage} (combined with
\Cref{v4-arith:w7-d:prop:density-breaks}) admits no elementary
reformulation.
\end{theorem}

\begin{proof}
Combine:
\Cref{v4-arith:w6-a:thm:HFD-ATP} (HFD closure),
\Cref{v4-arith:w7-d:thm:positive-cone} (Gaussian positive cone
closure),
\Cref{v4-arith:w7-d:prop:density-breaks} (density extension breaks),
and
\Cref{v4-arith:w6-a:thm:irreducible-ATP} (bounded-interval residual
RH-equivalent).
\end{proof}

\begin{remark}[route~7 D deliverable]
\label{v4-arith:w7-d:rem:deliverable}
Route~7 D produces three deliverables.

\begin{description}
\item[\textbf{(a) Gaussian sub-class closure.}] A-TP closes
unconditionally on the Gaussian sub-class for
$\lambda\geq\lambda_{c}$ with $\lambda_{c}\approx 2.52$
(\Cref{v4-arith:w7-d:thm:gauss-ATP}).

\item[\textbf{(b) Scalar characterisation.}] The Gaussian A-TP
range is exactly $[\lambda_{c},\infty)$
(\Cref{v4-arith:w7-d:cor:gauss-range}) with $\lambda_{c}$ the unique
root of the explicit erfc series $\widetilde{I}(\lambda)=0$
(\Cref{v4-arith:w7-d:thm:lambda-c}).

\item[\textbf{(c) Density extension analysis.}] Density-of-Gaussians
alone does not extend A-TP to full PW
(\Cref{v4-arith:w7-d:prop:density-breaks}); the quadratic structure
of the inequality requires positive-cone closure, not just linear
density. A-TP extends to the positive Gaussian cone
$\mathcal{C}_{\mathrm{Gauss}}^{+}$
(\Cref{v4-arith:w7-d:thm:positive-cone}) but no further via this
route.
\end{description}
\end{remark}

\section{Summary}
\label{v4-arith:w7-d:summary}

\begin{enumerate}
\item Gaussian self-duality and convolution closure collapse A-TP on
the Gaussian sub-class to a scalar inequality
$I(\lambda)\geq 0$ parameterised by a single $\lambda>0$
(\Cref{v4-arith:w7-d:cor:scalar-reduction}).

\item The scalar integral $I(\lambda)$ admits a closed-form
erfc series \eqref{v4-arith:w7-d:eq:I-erfc-series} via Hurwitz
expansion of the digamma kernel
(\Cref{v4-arith:w7-d:thm:erfc-series}); the series is absolutely
convergent and numerically tractable.

\item Asymptotic analysis: $I(\lambda)<0$ for $\lambda\to 0^{+}$
(narrow Gaussian samples $g(0)<0$;
\Cref{v4-arith:w7-d:prop:narrow}); $I(\lambda)>0$ for
$\lambda\to\infty$ (wide Gaussian samples the log-asymptotic region
of $g$; \Cref{v4-arith:w7-d:prop:wide}).

\item A unique critical $\lambda_{c}\approx 2.52$ exists with
$I(\lambda_{c})=0$; the sign of $I$ is monotone in $\lambda$
(\Cref{v4-arith:w7-d:thm:lambda-c},
\Cref{v4-arith:w7-d:prop:lambda-c-numerical}).

\item A-TP on the Gaussian sub-class closes unconditionally for
$\lambda\geq\lambda_{c}$
(\Cref{v4-arith:w7-d:thm:gauss-ATP},
\Cref{v4-arith:w7-d:cor:gauss-range}).

\item Extension by Gaussian density does not close full A-TP:
the quadratic structure of the inequality requires positive-cone
closure (not linear density)
(\Cref{v4-arith:w7-d:prop:density-breaks}).
A-TP extends unconditionally to the positive Gaussian cone
$\mathcal{C}_{\mathrm{Gauss}}^{+}$
(\Cref{v4-arith:w7-d:thm:positive-cone}).

\item The combined closed class
$\mathrm{PW}^{\mathrm{hol},\mathrm{HFD}}_{F}\cup\mathcal{C}_{\mathrm{Gauss}}^{+}$
strictly enlarges either the route~6 HFD class or the route~7 D
Gaussian cone alone; the complement in $\mathrm{PW}^{\mathrm{hol}}_{F}$
remains the irreducible RH-equivalent A-TP residual
(\Cref{v4-arith:w7-d:thm:wave7-d-irreducible}).
\end{enumerate}

\begin{remark}[Beilinson-principle closure]
\label{v4-arith:w7-d:rem:beilinson}
The route~7 D attack obeys the Beilinson principle: A-TP closes on
the specific Gaussian-positive-cone sub-class where the quadratic
structure is compatible with the sign-indefinite kernel $g$. The
critical parameter $\lambda_{c}$ is explicit and computable; the
admissible range $[\lambda_{c},\infty)$ is characterised sharply.
The elementary ceiling is reached: further progress requires either
(i) a non-positive-combination structural theory (e.g., de Branges
or Connes--Deninger), (ii) a non-Gaussian test class (e.g., compactly
supported PW), or (iii) a non-elementary spectral-density refinement
linking the zeros of $\zeta$ to the moment structure of Gaussian
integrals against $g$.
\end{remark}

\section{Bibliography}
\label{v4-arith:w7-d:bibliography}

\begin{description}
\item[Erd\'elyi--Magnus--Oberhettinger--Tricomi 1953.] \emph{Higher
Transcendental Functions} (Bateman Manuscript Project) Vol.~II,
Ch.~IX. McGraw--Hill. Erfc integrals, parabolic cylinder functions.

\item[Gradshteyn--Ryzhik 1965.] \emph{Tables of Integrals, Series,
and Products}. Academic Press. \S 3.323 Gaussian integrals, \S 4.335
log-weight Gaussian, \S 7.4 erfc tables, \S 8.251 complementary error
function identities.

\item[Titchmarsh 1948.] \emph{Introduction to the Theory of Fourier
Integrals} (2nd edn, Oxford). \S 2.3 density of Gaussians in the
Schwartz class; continuity of tempered-distribution functionals.

\item[Rudin 1973.] \emph{Functional Analysis}, McGraw--Hill.
Ch.~7 distributions, Paley--Wiener spaces, $L^{2}$ continuity
of tempered pairings.

\item[Hurwitz 1882.] Hurwitz, A. \emph{Einige Eigenschaften der
Dirichlet'schen Functionen.} Z.~Math.~Phys.~27. Hurwitz series
for the digamma function (cross-reference to the upstream analysis of
Ch~\ref{v4-arith:w6-a:ATP-direct-attack}).

\item[Abramowitz--Stegun 1964.] Abramowitz, M.; Stegun, I.A.
\emph{Handbook of Mathematical Functions}. National Bureau of
Standards Applied Math.~Series 55. \S 7 error functions (tables and
asymptotics); \S 6 gamma and digamma asymptotics.

\item[Weil 1952.] (Cross-reference.) Sur les formules explicites de
la th\'eorie des nombres premiers. Oeuvres Scientifiques Vol.~II.
Used as the Weil-explicit-formula input (upstream, not rederived).
\end{description}

\noindent\emph{Disjointness audit.} The primary catalogs for this
chapter (Erd\'elyi 1953 for erfc integral identities;
Gradshteyn--Ryzhik 1965 for tabulated Gaussian integrals;
Titchmarsh 1948 for Fourier integral continuity; Rudin 1973 for
distributional extension; Abramowitz--Stegun 1964 for erfc tables)
are classical analytic-tool catalogs that do not appear as derivation
sources in Chapters~\ref{v4-arith:closure}, \ref{v4-arith:kp-compat},
\ref{v4-arith:kms-gns-full}, \ref{v4-arith:w5-a:chapter},
\ref{v4-arith:w5-i:chapter}, or \ref{v4-arith:w6-a:ATP-direct-attack}.
Hurwitz 1882 and Weil 1952 appear upstream as derivation sources; the
present chapter uses them via cross-reference only (the Hurwitz series
is imported from Ch~\ref{v4-arith:w6-a:ATP-direct-attack}, and Weil's
explicit formula is imported from Ch~\ref{v4-arith:w5-a:chapter}).
The disjointness condition for this chapter passes by
per-source check.
